\begin{abstract}
Learned PID tuners are usually evaluated against weak baselines, on unvalidated models, with a handful of seeds.
We study whether a reinforcement-learning policy can improve the position PID gains of a geared brushed DC-motor joint beyond a \emph{validation-tuned constant} gain set, on an open, datasheet-derived digital twin that we describe as \prelim: it has not been compared with hardware.
Across \ConfThreeStaticN{} to \ConfThreeChangesN{} feasible scenarios per regime we compare five baselines, a hindsight oracle and PPO, SAC and TD3 policies, with the baseline gains and the reported policy seed chosen on separate validation seeds, return as the primary metric, one pre-declared primary test, and infeasible scenarios reported separately.
An unconstrained multiplicative PPO setup failed (median return difference to the constant $\VTwoStaticMed$, unstable training). A residual PPO acting around the constant improved on it by $\ConfThreeStaticMed$ (static), $\ConfThreeChangesMed$ (payload and target changes) and $\ConfThreeOodMed$ (out-of-distribution plants) median return on fresh confirmatory seeds, positive for all three training seeds, but it captured only $\ConfThreeStaticCapMean$\,\%, $\ConfThreeChangesCapMean$\,\% and $\ConfThreeOodCapMean$\,\% of the hindsight-oracle headroom on average over seeds. TD3 did not beat the constant and SAC was mixed.
All results are simulated on an unvalidated twin; the hardware part of the work is a replication protocol without results. Code, policies, per-episode data and the scripts that regenerate every table and figure are provided.
\end{abstract}
